% 第 13 题：含平行板电容器的电路与磁场区域（TikZ 线稿）
\begin{tikzpicture}[font=\small, >=Stealth, scale=1.0]
  % 左侧电路：E0、R1、R2
  \draw (0,0) to[battery1, l=$E_0$] (0,3);
  \draw (0,0) -- (1.2,0);
  \draw (1.2,-0.25) rectangle (2.8,0.25);
  \node[below] at (2.0,-0.25) {$R_1$};
  \draw (2.8,0) -- (3.7,0) -- (3.7,3) -- (0,3);
  \draw (3.7,3) -- (3.7,2.7);
  \draw (3.45,2.7) rectangle (3.95,1.0);
  \node[right] at (4.0,1.85) {$R_2$};
  \draw (3.7,1.0) -- (3.7,0);
  % 极板
  \draw[line width=1.8pt] (4.35,2.45) -- (9.55,2.45);
  \draw[line width=1.8pt] (4.35,0.45) -- (9.55,0.45);
  \draw (3.7,2.7) -- (4.35,2.45);
  \draw (3.7,1.0) -- (4.35,0.45);
  % 磁场（垂直纸面向里）
  \foreach \x in {5.1,5.9,6.7,7.5,8.3,9.1}{
    \foreach \y in {0.95,1.45,1.95}{
      \node at (\x,\y) {$\times$};
    }
  }
  % A 点与初速度
  \fill (4.42,0.45) circle (2pt) node[below left] {$A$};
  \draw[->, thick] (4.5,0.72) -- (5.55,0.72) node[midway, above] {$v$};
  % 板间距 d
  \draw[<->] (4.65,0.45) -- (4.65,2.45) node[midway, left] {$d$};
  % 板长
  \draw[<->] (4.35,0.10) -- (9.55,0.10) node[midway, below] {$\sqrt{3}d$};
  % 虚线
  \draw[dashed] (9.7,-0.25) -- (9.7,3.0);
  \node[right] at (9.78,2.45) {$E'$};
\end{tikzpicture}
